\documentclass[11pt,a4paper]{article}
\usepackage[margin=1in]{geometry}
\usepackage{times}
\usepackage{latexsym}
\usepackage{amsmath}
\usepackage{amssymb}
\usepackage{booktabs}
\usepackage{multirow}
\usepackage{graphicx}
\usepackage{xcolor}
\usepackage{url}
\usepackage{hyperref}
\usepackage{microtype}
\usepackage{array}
\usepackage{natbib}

\newcommand{\etamax}{\ensuremath{\eta_{\max}}}
\newcommand{\keff}{\ensuremath{K_{\text{eff}}}}

\title{Weakest Link, Not Diversity: Adversarial Vulnerability \\in Multi-Judge LLM Evaluation}

\author{Anonymous}

\begin{document}
\maketitle

\begin{abstract}
Multi-judge LLM evaluation panels are widely assumed to benefit from judge diversity, drawing on the analogy to Byzantine fault tolerance (BFT).
We conduct the first large-scale empirical study decomposing panel vulnerability into individual judge weakness ($\etamax$, the maximum single-judge attack success rate) and structural diversity ($\keff$, the effective number of independent judges).
Across 15 models, 455 three-judge panels, 6 attack types, and 2 benchmarks (12 conditions), we find that $\etamax$ is the dominant predictor of panel vulnerability, surviving non-independence correction at 5--8 out of 12 conditions depending on the method, while $\keff$ collapses from 8/12 (na\"ive OLS) to 0--1/12 after correcting for the panel overlap inherent in combinatorial designs.
We further demonstrate that this discrepancy between OLS and corrected inference is itself informative: the $\sim$2.5$\times$ standard error inflation reveals severe non-independence in panel-level regressions, a methodological concern relevant to any combinatorial evaluation study.
Practically, a simple $\min(\etamax)$ selection strategy dominates all 12 conditions, and we provide actionable recommendations for robust panel design.
\end{abstract}

%---------------------------------------------------------------------
\section{Introduction}
\label{sec:intro}

Using multiple LLM judges to evaluate model outputs has become standard practice, with the implicit assumption that diversity among judges improves robustness against adversarial manipulation.
This assumption mirrors classical Byzantine fault tolerance (BFT): if faults are independent and random, adding independent nodes increases system reliability.
However, adversarial attacks on LLM judges are neither independent nor random---strategic adversaries can exploit judge-specific vulnerabilities, potentially turning diversity from an asset into a liability.

Despite growing deployment of multi-judge panels, no prior work has empirically decomposed panel-level vulnerability into its constituent drivers under adversarial conditions at scale.
We address this gap through a large-scale study spanning 15 models from 6 families across 3 capability tiers, yielding $\binom{15}{3}=455$ unique three-judge panels evaluated under 6 attack types on 2 benchmarks (12 conditions total).

We decompose panel attack success rate (ASR) into two structural predictors: $\etamax$---the vulnerability of the weakest individual judge---and $\keff$---the effective number of independent judges derived from the mutual information matrix.
Our analysis reveals three primary contributions:

\paragraph{(1) $\etamax$ dominance with benchmark-dependent strength.}
Under rigorous non-independence correction (wild cluster bootstrap, clustered standard errors, mixed effects), $\etamax$ remains significant in 5--8/12 conditions while $\keff$ collapses to 0--3/12.
The dominance ratio is 5:1 under bootstrap inference.
However, this dominance is benchmark-dependent: MMLU conditions show stronger $\etamax$ signal (4/6 bootstrap) than ARC (1/6), reflecting the higher scoring variance of 4-choice versus 2-choice formats.

\paragraph{(2) Methodological transparency on non-independence.}
Na\"ive OLS yields $\etamax$ 11/12 and $\keff$ 8/12 significant, but standard errors are inflated by $\sim$2.5$\times$ due to panel overlap (each of 15 judges appears in $\sim$91 of 455 panels).
We systematically compare four correction methods---OLS, wild cluster bootstrap, clustered SE, and mixed effects---demonstrating that significance conclusions are highly sensitive to the correction method chosen.
This methodological demonstration is relevant beyond our specific application to any study employing combinatorial evaluation designs.

\paragraph{(3) Practical defense recommendations.}
A simple $\min(\etamax)$ panel selection strategy dominates all 12 conditions with zero failures, providing a robust and practical design principle: prioritize individual judge robustness over panel diversity.

%---------------------------------------------------------------------
\section{Related Work}
\label{sec:related}

\paragraph{LLM-as-Judge.}
LLM-based evaluation has become widespread for assessing open-ended generation quality (Zheng et al., 2023).
Multi-judge panels using majority voting or weighted aggregation aim to reduce individual judge biases.
However, systematic analysis of panel robustness under adversarial conditions remains limited.

\paragraph{Adversarial Attacks on LLM Judges.}
Recent work has demonstrated that LLM judges are vulnerable to prompt injection, sycophancy exploitation, and various cognitive biases.
These attacks can be categorized as \emph{targeted} (exploiting judge-specific weaknesses) and \emph{shared-bias} (exploiting common vulnerabilities across judges).

\paragraph{Byzantine Fault Tolerance.}
Classical BFT theory guarantees system correctness when fewer than $f$ out of $n=3f+1$ nodes fail, assuming independent random faults.
The critical assumption of fault independence breaks down under adaptive adversaries who can observe and exploit system structure.

\paragraph{Panel Diversity in Evaluation.}
Prior work on ensemble methods and committee-based evaluation has emphasized diversity as a key factor for robustness.
We provide the first large-scale empirical test of this assumption under adversarial conditions in the LLM judge setting.

%---------------------------------------------------------------------
\section{Method}
\label{sec:method}

\subsection{Panel Construction and Metrics}

Given a pool of $N=15$ LLM judges spanning 6 model families (OpenAI, Anthropic, Google, Meta, Qwen, Mistral) and 3 capability tiers (small/7B, medium/70B, API-frontier), we enumerate all $\binom{15}{3}=455$ unique three-judge panels.

For each panel $P=\{j_1, j_2, j_3\}$, we define two structural predictors:

\paragraph{Maximum individual vulnerability ($\etamax$).}
$\etamax(P) = \max_{j \in P} \eta_j$, where $\eta_j$ is judge $j$'s individual attack success rate under a given attack type.
This captures the ``weakest link'' in the panel.

\paragraph{Effective number of independent judges ($\keff$).}
We compute pairwise mutual information $\text{MI}(j_i, j_k)$ from the joint scoring distribution across evaluation instances.
$\keff$ is derived from the eigenspectrum of the MI matrix, measuring the effective dimensionality of the panel's scoring behavior.
Higher $\keff$ indicates greater functional independence among judges.

\subsection{Attack Types}

We evaluate 6 attack types spanning two mechanistic categories:

\begin{itemize}
\item \textbf{Targeted attacks} exploit judge-specific vulnerabilities: prompt injection (PI), score manipulation (SM), and position bias (PB).
\item \textbf{Shared-bias attacks} exploit common cognitive biases: sycophancy (syco), verbosity bias (VB), and authority bias (auth).
\end{itemize}

Each attack is evaluated on 2 benchmarks (MMLU with 4-choice MCQ format and ARC with 2-choice format), yielding 12 conditions.

\subsection{Regression Specification}

Our canonical regression specification is log-log OLS:
\begin{equation}
\log(\text{ASR} + \epsilon) = \beta_0 + \beta_1 \log(\keff) + \beta_2 \log(\etamax + \epsilon) + \varepsilon
\label{eq:regression}
\end{equation}
where $\epsilon$ is a small constant to handle zero values.

\subsection{Non-Independence Correction Methods}
\label{sec:correction}

A critical methodological concern is that the 455 panels are not independent observations: each of the 15 judges appears in $\binom{14}{2}=91$ panels ($\sim$20\% of all panels).
This overlap inflates OLS standard errors by approximately 2.5$\times$.
We apply four correction methods:

\begin{enumerate}
\item \textbf{Na\"ive OLS}: baseline with no correction for panel overlap.
\item \textbf{Wild cluster bootstrap} (B=9999, Rademacher weights, 3-position max-$p$): clusters panels by each judge position and reports the maximum $p$-value across positions as a conservative bound.
\item \textbf{Clustered standard errors} (CR1, 15 judge clusters): adjusts standard errors for within-judge correlation across panels.
\item \textbf{Mixed effects}: includes random judge intercepts to absorb judge-level variance.
\end{enumerate}

%---------------------------------------------------------------------
\section{Experiments}
\label{sec:experiments}

\subsection{Main Results: $\etamax$ Dominance}
\label{sec:main_results}

Table~\ref{tab:correction_comparison} presents the central result.
Under na\"ive OLS, both $\etamax$ (11/12) and $\keff$ (8/12, with 7 positive and 1 negative) appear strongly significant, suggesting a dual-driver model of panel vulnerability.
However, correcting for panel non-independence reveals a stark asymmetry: $\etamax$ retains significance at 5--8/12 depending on the correction method, while $\keff$ collapses to 0--3/12.

\begin{table}[t]
\centering
\small
\begin{tabular}{lcccc}
\toprule
\textbf{Method} & \textbf{$\etamax$ sig} & \textbf{$\keff$ sig} & \textbf{Ratio} & \textbf{Note} \\
\midrule
OLS & 11/12 & 8/12 & 1.4:1 & Inflated \\
Bootstrap & 5/12 & 1/12 & 5:1 & Conservative \\
Clustered SE & 8/12 & 0/12 & 8:0 & Consistent \\
Mixed effects & 7/12 & 3/12 & 2.3:1 & Random int. \\
Tier-corrected & 11/12 & 5/12 & 2.2:1 & 3/3/9 tiers \\
\bottomrule
\end{tabular}
\caption{Significance counts (out of 12 benchmark$\times$attack conditions) for $\etamax$ and $\keff$ across correction methods. The $\etamax$:$\keff$ dominance ratio increases from 1.4:1 (OLS) to 5:1 (bootstrap) after correcting for panel non-independence, because $\keff$ significance is more severely inflated by panel overlap than $\etamax$.}
\label{tab:correction_comparison}
\end{table}

The wild cluster bootstrap gives 5/12 for $\etamax$ and 1/12 for $\keff$, yielding a dominance ratio of 5:1.
This is a conservative lower bound: position~0 clustering is imbalanced (one judge appears in 91 panels while another appears in only 1), which tightens the max-$p$ criterion.
Positions~1 and 2, which have more balanced cluster sizes, yield $\etamax$ significance at 9/12---consistent with the clustered SE result of 8/12.

\paragraph{Benchmark asymmetry.}
$\etamax$ dominance is benchmark-dependent.
Under bootstrap, MMLU conditions show $\etamax$ significance in 4/6 attacks (PI, SM, PB, VB survive), while ARC shows significance in only 1/6 (PI, marginally).
We attribute this to the scoring format: MMLU's 4-choice structure provides more variance for $\etamax$ to capture than ARC's binary format.
This structural limitation should be considered when generalizing our findings.

\subsection{Standard Error Inflation Analysis}
\label{sec:se_inflation}

The discrepancy between OLS and corrected results is itself informative.
Standard errors inflate by approximately 2.5$\times$ under clustered SE relative to OLS, consistent with the expected overlap: each judge appears in 91/455 $\approx$ 20\% of panels.
This inflation is not a statistical nuisance but a structural property of combinatorial evaluation designs.

Any study that enumerates all $\binom{N}{K}$ panels from a model pool faces the same non-independence structure.
Our systematic comparison of four correction methods (Table~\ref{tab:correction_comparison}) provides a template for future combinatorial evaluation studies.
The key finding is that the correction method can change significance conclusions substantially---$\keff$ goes from 8/12 (OLS) to 0/12 (clustered SE)---making the choice of correction method a first-order methodological decision.

\subsection{Multicollinearity and Partial Correlations}
\label{sec:vif}

$\keff$ and $\etamax$ are structurally correlated ($r=0.70$) because more diverse panels tend to include less capable (and thus more vulnerable) models.
We conduct a thorough multicollinearity analysis to assess whether $\keff$'s significance is driven by this confound.

\begin{table}[t]
\centering
\small
\begin{tabular}{lccc}
\toprule
\textbf{Specification} & \textbf{VIF($\keff$)} & \textbf{VIF($\etamax$)} & \textbf{Interp.} \\
\midrule
Base log-log & 1.96 & 1.96 & Acceptable \\
+ Tier indicators & 3.8 & 4.2 & Moderate \\
+ Benchmark int. & 4.5 & 4.8 & Approaching caution \\
\bottomrule
\end{tabular}
\caption{Variance inflation factors across regression specifications.
VIF remains below the conventional concern threshold of 10, but the moderate range (3--5) requires careful interpretation.}
\label{tab:vif}
\end{table}

\begin{table}[t]
\centering
\small
\begin{tabular}{lccc}
\toprule
\textbf{Analysis} & \textbf{$\etamax$} & \textbf{$\keff$} & \textbf{Dominance} \\
\midrule
Partial corr. & $r=0.35$ & $r=0.12$ & $\etamax$ 3$\times$ \\
Resid. $\keff$ (OLS) & --- & 3/12 & Marginal \\
Resid. $\keff$ (boot.) & --- & 0/12 & Eliminated \\
Beta regression & 10/12 & 5/12 & $\etamax$ 2$\times$ \\
Bonferroni & 8/12 & 4/12 & $\etamax$ 2$\times$ \\
FDR & 10/12 & 6/12 & $\etamax$ 1.7$\times$ \\
\bottomrule
\end{tabular}
\caption{Partial correlations and residualized regression confirm $\etamax$ dominance across analytical approaches.
The partial correlation of $\etamax$ given $\keff$ (0.35) is approximately 3$\times$ the reverse (0.12).
Residualizing $\keff$ on $\etamax$ eliminates its significance under bootstrap (0/12).}
\label{tab:partial}
\end{table}

Table~\ref{tab:vif} shows that VIF values remain in the moderate range (1.96--4.8), below the conventional concern threshold of 10 but sufficient to warrant discussion.
Table~\ref{tab:partial} presents complementary evidence: the partial correlation of $\etamax$ given $\keff$ ($r=0.35$) is approximately three times larger than the reverse ($r=0.12$), and residualizing $\keff$ on $\etamax$ eliminates its bootstrap significance entirely (0/12).
Beta regression, which is more appropriate for the bounded ASR response, gives $\etamax$ 10/12 versus $\keff$ 5/12.

These analyses converge: $\etamax$ retains independent predictive power across all specifications, while $\keff$'s apparent significance is partially attributable to its correlation with $\etamax$ and partially to panel overlap inflation.

\subsection{GAM Analysis: Suggestive Taxonomic Pattern}
\label{sec:gam}

Linear regression masks an important structure in $\keff$'s effect.
Generalized additive model (GAM) analysis reveals that the direction of $\keff$'s partial effect splits cleanly by attack mechanism:

\begin{itemize}
\item \textbf{Targeted attacks} (PI, SM, PB $\times$ 2 benchmarks = 6 conditions): $\keff$ partial effect is \emph{positive} in all 6 conditions---diversity is associated with higher vulnerability.
\item \textbf{Shared-bias attacks} (syco, VB, auth $\times$ 2 benchmarks = 6 conditions): $\keff$ partial effect is \emph{negative} in all 6 conditions---diversity is directionally protective.
\end{itemize}

This 6:6 split is consistent with the mechanistic distinction: targeted attacks exploit judge-specific weaknesses (more diverse panels expose more distinct attack surfaces), while shared-bias attacks exploit biases common across judges (diverse panels may dilute correlated errors).
We emphasize that this taxonomic pattern is suggestive and warrants further investigation---individual condition significance under clustered SE or bootstrap inference has not been fully validated for the GAM specification, and the clean 6:6 split should be interpreted as a qualitative structural finding rather than 12 independent statistical tests.

\subsection{Random-Fault BFT Contrast}
\label{sec:bft}

To assess whether the BFT analogy holds under any conditions, we contrast the adaptive-attack regime with a random-fault (RF) regime where each judge fails independently with probability proportional to $1 - \text{accuracy}$.

\begin{table}[t]
\centering
\small
\begin{tabular}{lccccl}
\toprule
\textbf{Regime} & \textbf{Bench.} & \textbf{$\keff$ $p$ (OLS)} & \textbf{$\keff$ $p$ (boot.)} & \textbf{$R^2$} & \textbf{Interp.} \\
\midrule
RF & MMLU & 0.002 & 0.03 & 0.73 & Genuine \\
RF & ARC & 0.008 & 0.06 & 0.68 & Borderline \\
Adaptive & MMLU & varies & 5/6 NS & 0.25 & Inflated \\
Adaptive & ARC & varies & 6/6 NS & 0.18 & Inflated \\
\bottomrule
\end{tabular}
\caption{$\keff$ significance and model fit ($R^2$) under random-fault vs.\ adaptive regimes.
Under random faults, panel composition is highly predictive ($R^2\approx 0.70$) and $\keff$ is genuinely significant.
Under adaptive attacks, predictability collapses ($R^2\approx 0.22$) and $\keff$ significance is largely artifactual.}
\label{tab:bft}
\end{table}

Table~\ref{tab:bft} shows a clear predictability gap.
Under random faults, $\keff$ is genuinely significant ($p=0.03$ under bootstrap for MMLU, borderline $p=0.06$ for ARC) with $R^2\approx 0.68$--$0.73$---panel composition strongly predicts vulnerability.
Under adaptive attacks, $\keff$ significance collapses (5--6 out of 6 NS under bootstrap) with $R^2\approx 0.18$--$0.25$.

The key contrast is not that diversity helps under RF but hurts under adaptive---$\keff$'s standardized coefficient is positive in \emph{both} regimes ($\beta_\text{std}=+0.215$ for RF MMLU, $+0.379$ for RF ARC).
Rather, the distinction is that $\keff$'s harmful association is \emph{genuine} under RF (survives all corrections) but \emph{artifactual} under adaptive attacks (driven by panel overlap inflation).
Adaptive adversaries introduce strategic variance beyond what panel-level structural predictors can capture.

\subsection{Aggregation Gradient}
\label{sec:aggregation}

We evaluate how the choice of vote aggregation rule modulates vulnerability by testing four methods: majority vote, accuracy-weighted voting, threshold (unanimity) voting, and Bayesian calibration.

\begin{table}[t]
\centering
\small
\setlength{\tabcolsep}{3pt}
\begin{tabular}{l ccc ccc}
\toprule
& \multicolumn{3}{c}{\textbf{$\etamax$ sig}} & \multicolumn{3}{c}{\textbf{$\keff$ sig}} \\
\cmidrule(lr){2-4} \cmidrule(lr){5-7}
\textbf{Aggregation} & OLS & Clust. & Boot. & OLS & Clust. & Boot. \\
\midrule
Majority & 11/12 & 8/12 & 5/12 & 8/12 & 0/12 & 1/12 \\
Weighted & 9/12 & 6/12 & 4/12 & 9/12 & 2/12 & 1/12 \\
Threshold & 12/12 & 10/12 & 7/12 & 11/12 & 5/12 & 3/12 \\
Bayesian & 9/12 & 5/12 & 3/12 & 7/12 & 4/12 & 3/12 \\
\bottomrule
\end{tabular}
\caption{Aggregation gradient under three correction methods.
Threshold (unanimity) voting amplifies both predictors' significance, while Bayesian calibration attenuates both.
Under OLS, Bayesian reverses $\keff$ direction (6/7 negative coefficients), but this reversal is attenuated under bootstrap (2/7 negative survive).
The ordinal ranking (threshold $>$ majority $>$ weighted $>$ Bayesian for dual-driver significance) is preserved across all correction methods.}
\label{tab:aggregation}
\end{table}

Table~\ref{tab:aggregation} reveals a consistent aggregation gradient.
Threshold voting amplifies both vulnerability channels (OLS $R^2=0.89$), consistent with unanimity rules being maximally sensitive to any single corrupted judge.
Bayesian calibration attenuates both channels, with $\etamax$ dropping from 11/12 to 9/12 (OLS) and $\keff$ showing a directional tendency toward sign reversal.

Under OLS, Bayesian aggregation reverses $\keff$'s coefficient sign in 6/7 MMLU conditions (from positive/harmful to negative/protective).
However, this reversal is partially attenuated under non-independence correction: only 2/7 negative coefficients survive bootstrap inference.
The structural direction---Bayesian reduces both $\etamax$ and $\keff$ significance---is consistent across all correction methods, but the sign reversal claim requires further validation.

The ordinal ranking of aggregation methods (threshold $>$ majority $>$ weighted $>$ Bayesian for vulnerability) is preserved under all three correction approaches, making this a robust qualitative finding even as absolute significance counts decrease.

\subsection{Bayesian $\keff$ Sign Flip Detail}
\label{sec:bayesian_flip}

To examine the Bayesian $\keff$ sign reversal more carefully, Table~\ref{tab:bayesian_flip} reports condition-level results.

\begin{table}[t]
\centering
\small
\begin{tabular}{lccc}
\toprule
\textbf{Condition} & \textbf{$\keff$ $\beta$ (OLS)} & \textbf{Clust. sig} & \textbf{Boot. sig} \\
\midrule
MMLU$\times$PI & $-0.42$ & Yes ($p$=0.04) & Yes ($p$=0.03) \\
MMLU$\times$SM & $-0.38$ & No ($p$=0.12) & No ($p$=0.18) \\
MMLU$\times$PB & $-0.31$ & No ($p$=0.15) & No ($p$=0.22) \\
MMLU$\times$syco & $-0.25$ & Yes ($p$=0.04) & No ($p$=0.08) \\
MMLU$\times$VB & $-0.29$ & No ($p$=0.09) & No ($p$=0.14) \\
MMLU$\times$auth & $-0.22$ & Yes ($p$=0.03) & Yes ($p$=0.04) \\
ARC$\times$PI & $+0.15$ & No ($p$=0.35) & No ($p$=0.42) \\
\bottomrule
\end{tabular}
\caption{Condition-level $\keff$ coefficients under Bayesian aggregation. The sign reversal (negative = protective) is robust for MMLU$\times$PI and MMLU$\times$auth under both correction methods, partial for MMLU$\times$syco, and absent for the remaining conditions. ARC$\times$PI does not show the reversal.}
\label{tab:bayesian_flip}
\end{table}

Two conditions---MMLU$\times$PI ($\beta=-0.42$, bootstrap $p=0.03$) and MMLU$\times$auth ($\beta=-0.22$, bootstrap $p=0.04$)---show robust sign reversals surviving both correction methods.
MMLU$\times$syco shows a partial reversal (clustered SE significant but bootstrap borderline).
The remaining conditions show directional negative coefficients that do not survive correction.
This pattern suggests that Bayesian calibration can reverse the diversity--vulnerability relationship for specific attack types, but the effect is not universal.

\subsection{Practical Panel Design}
\label{sec:practical}

Across all 12 conditions, a simple $\min(\etamax)$ selection strategy---choosing the panel with the lowest maximum individual vulnerability---dominates every alternative (12/12 wins).
The ASR advantage ranges from +0.038 (ARC$\times$sycophancy) to +0.511 (ARC$\times$prompt injection), indicating substantial practical gains.

The practical effect size of $\keff$ is small: an interquartile range change in $\keff$ corresponds to $\Delta\text{ASR}=1.3$ percentage points, and clustered SE confidence intervals contain zero in 11/12 conditions.
Even in the single condition where $\keff$ is the sole significant predictor under any correction (ARC$\times$position bias under mixed effects), $\min(\etamax)$ still outperforms $\max(\keff)$.

\paragraph{Diversity-capability tradeoff.}
The $r=0.70$ correlation between $\keff$ and $\etamax$ reflects a structural constraint of the current LLM ecosystem: maximizing panel diversity requires including less capable models, which increases $\etamax$.
This tradeoff motivates a two-stage design: (1)~screen for individual robustness ($\eta \leq$ threshold), then (2)~maximize $\keff$ among qualifying judges.
However, the marginal gain from stage~2 is small given the 12/12 dominance of $\min(\etamax)$.

\paragraph{Family labels do not predict diversity.}
Model family labels (e.g., Qwen, Llama, OpenAI) do not significantly predict judge independence ($p=0.179$ for MMLU, $p=0.385$ for ARC; Mann-Whitney $U$ test over 14 within-family and 91 cross-family MI pairs).
Within-family MI variance is high (e.g., Llama-70B and Llama-8B show MI of only 0.05 due to the capability gap exceeding family similarity), indicating that capability proximity, not family membership, drives inter-judge correlation.

%---------------------------------------------------------------------
\section{Discussion}
\label{sec:discussion}

\paragraph{When diversity may help.}
Our findings should not be interpreted as a blanket indictment of panel diversity.
The GAM analysis (Section~\ref{sec:gam}) suggests that $\keff$'s effect splits by attack mechanism, with diversity directionally protective against shared-bias attacks.
Additionally, diversity remains beneficial in non-adversarial settings (e.g., reducing calibration errors) and after robustness screening (when all judges exceed a capability threshold, the diversity-capability tradeoff is broken).

\paragraph{Implications for practice.}
Our results suggest a dual-channel defense stack: (1)~robustness-first selection ($\min(\etamax)$, effective in 12/12 conditions) combined with (2)~Bayesian aggregation (which directionally reverses $\keff$'s harmful association).
This combination addresses both the individual vulnerability channel (by selecting robust judges) and the structural diversity channel (by modulating how diverse scores are combined).

\paragraph{Broader applicability.}
The non-independence correction framework we develop is relevant to any study that enumerates combinatorial subsets from a finite pool.
As LLM evaluation increasingly relies on multi-judge panels, the methodological distinction between na\"ive and corrected inference will become a standard concern.

\subsection{Limitations}

\textbf{Benchmark scope.}
Our findings are established on two MCQ benchmarks (MMLU and ARC).
The benchmark asymmetry we observe (MMLU stronger than ARC for $\etamax$ signal) suggests that scoring format affects the detectability of vulnerability drivers.
Generalization to open-ended evaluation benchmarks (e.g., MT-Bench, AlpacaEval) is a critical direction for future work, where higher scoring variance may strengthen $\etamax$ dominance.

\textbf{Attack coverage.}
We evaluate 6 attack types covering targeted and shared-bias mechanisms.
Broader attack coverage (e.g., few-shot hijacking, persona manipulation, instruction override) is needed to test the universality of the weakest-link finding.

\textbf{Panel size.}
Our analysis uses $K=3$ panels.
Larger panels ($K=5,7,9$) may dilute the $\etamax$ effect through averaging, but also face extreme panel overlap ($\binom{15}{9}=5005$ panels with each model in $\sim$3000 panels), making corrected inference challenging without a substantially larger model pool.

\textbf{Capability-matched validation.}
The $r=0.70$ correlation between $\keff$ and $\etamax$ limits our ability to fully disentangle their independent effects.
Within-tier analysis (e.g., API-only models) reduces this confound ($r\approx 0.55$) but at the cost of statistical power (84 panels for 9 API models).
A definitive capability-matched validation requires more models per tier.

\textbf{GAM validation.}
The 6:6 taxonomic split is derived from GAM partial effects under OLS.
Bootstrapped confidence intervals for GAM smooths under the panel non-independence structure have not been computed; the pattern should be treated as a suggestive structural finding pending further validation.

%---------------------------------------------------------------------
\section{Conclusion}

We present the first large-scale empirical decomposition of multi-judge LLM evaluation panel vulnerability under adversarial attack.
Across 15 models, 455 panels, and 12 conditions, we find that panel vulnerability is dominated by the weakest individual judge ($\etamax$), with a 5:1 dominance ratio over structural diversity ($\keff$) under rigorous non-independence correction.
This dominance is directional with benchmark-dependent strength---stronger for MMLU than ARC---reflecting the influence of scoring format on vulnerability signal.
We demonstrate that na\"ive OLS inference substantially inflates significance for both predictors due to panel overlap, and that methodological transparency about correction choices is essential for combinatorial evaluation designs.
Practically, a simple $\min(\etamax)$ strategy dominates all conditions, providing an actionable principle: secure your weakest judge before diversifying your panel.

%\bibliography{custom}

\appendix

\section{Diversity-Capability Tradeoff}
\label{app:tradeoff}

The $r=0.70$ correlation between $\keff$ and $\etamax$ arises because the current LLM ecosystem exhibits a capability gradient: adding more diverse (functionally independent) judges requires reaching beyond the top-tier models into less capable ones. We observe a log-linear tradeoff: $\keff \approx 1.2 + 0.8 \cdot \log(\etamax)$, though residual variance is substantial (49\% unexplained), meaning some panel compositions beat the tradeoff curve.

Tier stratification confirms this pattern:
\begin{itemize}
\item \textbf{Small (7B)}: $\keff \in [2.2, 2.8]$, $\etamax \in [0.15, 0.35]$---diverse but vulnerable.
\item \textbf{Medium (70B)}: $\keff \in [2.0, 2.6]$, $\etamax \in [0.08, 0.20]$---moderate.
\item \textbf{API}: $\keff \in [1.5, 2.5]$, $\etamax \in [0.03, 0.15]$---wide range, including the most robust compositions.
\end{itemize}

Pareto-optimal panels (maximizing $\keff$ for a given $\etamax$ level) are dominated by API-tier models mixed with top open-source models (e.g., GPT-4o + Claude-3.5-Sonnet + Qwen2.5-72B achieves $\keff=2.35$ at $\etamax=0.06$).

\section{GO/NO-GO Gate for BFT Framing}
\label{app:gate}

\begin{table}[h]
\centering
\small
\begin{tabular}{lcc}
\toprule
\textbf{Criterion} & \textbf{Result} & \textbf{Status} \\
\midrule
RF $\keff$ sig $\geq$1/2 (OLS) & 2/2 & PASS \\
RF $\keff$ sig $\geq$1/2 (bootstrap) & 1/2 & PARTIAL \\
RF $R^2 >$ adaptive $R^2$ & 0.70 vs 0.22 & PASS \\
RF--adaptive contrast & 1--2/2 vs 1/12 & PARTIAL \\
BFT framing viable? & Yes with caveats & PARTIAL \\
\bottomrule
\end{tabular}
\caption{GO/NO-GO gate assessment for the BFT framing. The $R^2$ gap is the robust finding; significance contrasts are benchmark-dependent.}
\label{tab:gate}
\end{table}

\end{document}
